\chapter{Python API Reference}
\label{ch:api}

Complete reference for every public class and function of \qv{}~0.6.1.
Anything documented here is importable directly from the top-level package:
\code{from qanneal import ...}.

% ===========================================================================
\section{The high-level solver: \code{solve()}}
\label{sec:api-solve}

\begin{apibox}{solve(problem, method="sqa", **kwargs) $\rightarrow$ SolveResult}
One-call interface: accepts any supported problem type, auto-selects and
tunes a schedule, runs \code{reads} independent runs, returns the best
result.
\end{apibox}

\subsection{Accepted problem types}

The \code{problem} argument is auto-detected:

\begin{center}
\small
\begin{tabular}{@{}lp{8.4cm}@{}}
\toprule
\textbf{Type} & \textbf{Interpretation}\\
\midrule
\code{DenseIsing}, \code{SparseIsing} & used directly\\
\code{QUBO} & converted via \code{to\_ising()}\\
\code{np.ndarray} $(n\times n)$ & QUBO matrix, Eq.~\eqref{eq:qubo}\\
\code{dict\{(i,j): float\}} & sparse QUBO dictionary\\
\code{list[(i, j, float)]} & sparse QUBO entry list\\
\code{dimod.BinaryQuadraticModel} & BINARY or SPIN vartype (requires \code{dimod})\\
\code{networkx.Graph} & node attr \code{bias}, edge attr \code{weight} (requires \code{networkx})\\
\bottomrule
\end{tabular}
\end{center}

For dict/list inputs, pass \code{n=...} if the variable count cannot be
inferred. For BQM/graph inputs, \code{result.var\_order} maps positions back
to your original variable labels.

\subsection{Common parameters (all methods)}

\begin{center}
\small
\begin{tabular}{@{}llp{7.9cm}@{}}
\toprule
\textbf{Parameter} & \textbf{Default} & \textbf{Meaning}\\
\midrule
\code{method} & \code{"sqa"} & \code{"sa"}, \code{"sqa"}, \code{"sqapt"}, \code{"ctpimc"}.\\
\code{reads} & 1 & Independent runs; best over all reads is returned.\\
\code{sweeps\_per\_beta} & 20 & Metropolis sweeps per schedule step (also serves as \code{sweeps\_per\_step} in optimal mode).\\
\code{schedule} & \code{None} & Schedule object; auto-tuned per method if \code{None}.\\
\code{seed} & \code{None} & RNG seed; \code{None} = random.\\
\code{progress} & \code{True} & tqdm progress bar.\\
\code{backend} & \code{"cpu"} & Compute backend.\\
\code{n} & \code{None} & Size hint for dict/list inputs.\\
\code{return\_bits} & \code{False} & Return $\{0,1\}$ bits instead of $\{-1,+1\}$ spins.\\
\bottomrule
\end{tabular}
\end{center}

\subsection{SQA / SQAPT parameters}

\begin{center}
\small
\begin{tabular}{@{}llp{7.9cm}@{}}
\toprule
\textbf{Parameter} & \textbf{Default} & \textbf{Meaning}\\
\midrule
\code{trotter\_slices} & 32 & Imaginary-time slices $M$ (Section~\ref{sec:trotter}).\\
\code{replicas} & 1 & Parallel Trotter systems (\code{sqa}) or ladder rungs (\code{sqapt}).\\
\code{worldline\_sweeps} & 5 & Whole-worldline flips per step (Eq.~\eqref{eq:worldline}).\\
\code{cluster\_sweeps} & 0 & Swendsen--Wang segment flips per step; $>0$ enables +SW.\\
\code{continuous\_time\_slices} & 0 & If $>0$, overrides $M$ to approximate continuous time.\\
\bottomrule
\end{tabular}
\end{center}

\subsection{SQAPT-only parameters}

\begin{center}
\small
\begin{tabular}{@{}llp{7.9cm}@{}}
\toprule
\textbf{Parameter} & \textbf{Default} & \textbf{Meaning}\\
\midrule
\code{pt\_steps} & 50 & Local-update + swap epochs (standard schedule).\\
\code{swap\_interval} & 1 & Attempt swaps every $N$ epochs.\\
\code{pt\_betas} & \code{None} & Explicit $\beta$ ladder (overrides schedule).\\
\code{pt\_gammas} & \code{None} & Explicit $\Gamma$ ladder (must pair with \code{pt\_betas}).\\
\bottomrule
\end{tabular}
\end{center}

\subsection{CT-PIMC-only parameters}

\begin{center}
\small
\begin{tabular}{@{}llp{7.9cm}@{}}
\toprule
\textbf{Parameter} & \textbf{Default} & \textbf{Meaning}\\
\midrule
\code{ctpimc\_qubits\_per\_update} & 1 & Spins per cluster proposal.\\
\code{ctpimc\_qubits\_per\_chain} & 1 & Chain length for lattice-mode updates.\\
\bottomrule
\end{tabular}
\end{center}

\subsection{Optimal-schedule parameters (\code{schedule\_type="optimal"})}

Applicable to \code{sqa} and \code{sqapt}; theory in
Chapter~\ref{ch:optimal}.

\begin{center}
\small
\begin{tabular}{@{}llp{7.6cm}@{}}
\toprule
\textbf{Parameter} & \textbf{Default} & \textbf{Meaning}\\
\midrule
\code{schedule\_type} & \code{"standard"} & \code{"optimal"} activates the adaptive $\Jp$ schedule.\\
\code{optimal\_eps\_tilde} & 0.0 & $\le0$ = per-instance budget calibration (recommended); $>0$ = legacy fixed online scale.\\
\code{optimal\_alpha} & $15/14$ & Universality exponent $\alpha$ (Eq.~\eqref{eq:update-law}).\\
\code{optimal\_num\_steps} & 100 & Adaptive step budget $T$.\\
\code{optimal\_j\_perp\_start} & auto & $J_0$; from the tuned schedule's quantum end.\\
\code{optimal\_j\_perp\_end} & auto & $J_{\mathrm{end}}$; default $\max(J_0,\,5J_{\mathrm{rms}})$.\\
\code{optimal\_beta} & auto & Fixed $\beta$; SQA default $\max(4/J_{\mathrm{rms}},1)$.\\
\code{optimal\_beta\_ramp\_fraction} & 0.3 & SQA two-phase split (Section~\ref{sec:optimal-params}).\\
\code{optimal\_debug\_csv} & \code{""} & Path for the per-step diagnostic CSV.\\
\bottomrule
\end{tabular}
\end{center}

\subsection{\code{SolveResult}}

\begin{center}
\small
\begin{tabular}{@{}llp{8.2cm}@{}}
\toprule
\textbf{Field} & \textbf{Type} & \textbf{Meaning}\\
\midrule
\code{method} & \code{str} & Method used.\\
\code{samples} & \code{list[np.ndarray]} & Best spin/bit vector from each read.\\
\code{energies} & \code{list[float]} & Best energy from each read.\\
\code{best\_sample} & \code{np.ndarray} & Global best configuration.\\
\code{best\_energy} & \code{float} & Global best energy.\\
\code{trace} & \code{list[float]} & Energy trace from the first read.\\
\code{var\_order} & \code{list} & Variable ordering (BQM/graph inputs).\\
\code{return\_bits} & \code{bool} & Whether samples are bits or spins.\\
\bottomrule
\end{tabular}
\end{center}

Manual spin$\to$bit conversion:
\code{bits = ((result.best\_sample + 1) // 2).astype(int)}.

% ===========================================================================
\section{Model classes}
\label{sec:api-models}

\subsection{\code{DenseIsing}}

\begin{apibox}{DenseIsing(h, J, c=0.0)}
Fully-connected Ising model, dense $n\times n$ storage. Energy:
Eq.~\eqref{eq:ising}.
\end{apibox}

\begin{center}
\small
\begin{tabular}{@{}llp{7.7cm}@{}}
\toprule
\textbf{Argument} & \textbf{Type} & \textbf{Meaning}\\
\midrule
\code{h} & \code{ndarray (n,)} & Local fields; $h_i>0$ pushes spin $i$ toward $-1$.\\
\code{J} & \code{ndarray (n,n)} & Symmetric coupling matrix; upper triangle used; ferromagnetic $<0$, antiferromagnetic $>0$.\\
\code{c} & \code{float} & Constant offset (energy bookkeeping only).\\
\bottomrule
\end{tabular}
\end{center}

Methods: \code{size()} $\to$ \code{int}; \code{energy(spins)} $\to$
\code{float} ($O(n^2)$); \code{delta\_energy(spins, flip)} $\to$
\code{float} ($O(n)$, Eq.~\eqref{eq:deltaE}); \code{fields()} and
\code{couplings()} return the stored $h$ and $J$.

\subsection{\code{SparseIsing} and \code{SparseEdge}}

\begin{apibox}{SparseIsing(h, edges, n, c=0.0)}
Edge-list Ising model; scales to $n\geq10^5$. \code{energy()} is
$O(|E|)$, \code{delta\_energy()} is $O(\deg i)$.
\end{apibox}

\code{edges} is a \code{list[SparseEdge]}; \code{SparseEdge(i, j, value)}
holds one coupling $J_{ij}$ (one direction suffices, $i<j$; the model
symmetrises).

\begin{pycode}
from qanneal import SparseIsing, SparseEdge
import numpy as np

n = 5
edges = [SparseEdge(0, 1, 1.0), SparseEdge(1, 2, -0.5),
         SparseEdge(2, 3, 0.8), SparseEdge(3, 4, 1.2)]
ising = SparseIsing(np.zeros(n), edges, n)
\end{pycode}

\subsection{\code{QUBO}}

\begin{apibox}{QUBO(Q) \;|\; QUBO(entries, n) \;|\; QUBO(bqm)}
QUBO model, Eq.~\eqref{eq:qubo}. Constructors: dense
\code{ndarray}; sparse \code{list[(i,j,val)]} or
\code{dict\{(i,j): val\}} with explicit \code{n}; or a \code{dimod}
BQM (BINARY or SPIN, auto-converted).
\end{apibox}

Methods: \code{size()}; \code{to\_ising()} $\to$ \code{DenseIsing},
implementing the exact conversion Eq.~\eqref{eq:qubo2ising} (symmetrises
$Q$, applies $x=(1+s)/2$, carries the offset $c$).

% ===========================================================================
\section{Schedule classes and helpers}
\label{sec:api-schedules}

\subsection{\code{AnnealSchedule}}

\begin{apibox}{AnnealSchedule.linear(beta\_start, beta\_end, steps)}
Classical schedule: linear ramp of inverse temperatures. Attribute
\code{betas: list[float]}.
\end{apibox}

\subsection{\code{SQASchedule}}

\begin{apibox}{SQASchedule.from\_vectors(betas, gammas)}
Quantum schedule: paired $(\beta,\Gamma)$ sequences (equal lengths).
Attributes \code{betas}, \code{gammas}. $\Gamma$ should start high and
decay geometrically (see the warning in Chapter~\ref{ch:sqa}).
\end{apibox}

\subsection{Fixed helpers}

\begin{itemize}
\item \code{auto\_schedule\_sa(steps=50, beta\_start=0.1,
beta\_end=4.0)} $\to$ \code{AnnealSchedule} --- fixed linear ramp.
\item \code{auto\_schedule\_sqa(steps=50, beta\_start=0.1,
beta\_end=4.0, gamma\_start=5.0, gamma\_end=0.01)} $\to$
\code{SQASchedule} --- fixed ramp with geometric $\Gamma$ decay.
\end{itemize}

\subsection{Tuned helpers (recommended)}

Mechanics in Sections~\ref{sec:tuned-schedules} and \ref{sec:sqa-tuned}.

\begin{apibox}{auto\_schedule\_sa\_tuned(problem, mode="balanced", steps=None,
probes=256, seed=1234, n=None, beta\_end\_scale=1.0)}
Problem-adaptive SA schedule from the probed 75th-percentile
$|\Delta E|$ and acceptance targets, Eq.~\eqref{eq:beta-from-p}.
\end{apibox}

\begin{apibox}{auto\_schedule\_sqa\_tuned(..., gamma\_end\_scale=1.0)}
As above, additionally calibrating the $\Gamma$ ramp
($\Gamma_{\text{start}}=m_\Gamma\Delta$, geometric decay). Set
\code{gamma\_end\_scale=0.04} when the schedule feeds CT-PIMC.
\end{apibox}

\begin{apibox}{auto\_ladder\_sqa\_tuned(problem, replicas=8, mode="balanced",
probes=256, seed=1234, n=None)}
$(\beta,\Gamma)$ ladder for SQAPT: geometric $\beta$ rungs, inverse-geometric
$\Gamma$ rungs. Returns an \code{SQASchedule} whose entries are ladder
points, not a time sequence.
\end{apibox}

\subsection{Optimal-schedule helpers}

\begin{apibox}{j\_perp\_from\_beta\_gamma(beta, gamma, trotter\_slices=32)
$\rightarrow$ float}
The Trotter coupling of Eq.~\eqref{eq:jperp}:
$\Jp=\tfrac12\ln\bigl(1/\tanh(\beta\Gamma/M)\bigr)$.
\end{apibox}

\begin{apibox}{optimal\_j\_perp\_params(problem, mode="balanced",
trotter\_slices=32, probes=256, seed=1234, n=None)
$\rightarrow$ (beta, j\_perp\_start, j\_perp\_end)}
Problem-adaptive endpoints for \code{run\_optimal}: derives
$(\beta,\Gamma_{\text{start}},\Gamma_{\text{end}})$ from the tuned-schedule
probe and maps through Eq.~\eqref{eq:jperp}; the returned
$\beta=\max(4/J_{\mathrm{rms}},1)$.
\end{apibox}

\begin{apibox}{j\_rms\_from\_problem(problem) $\rightarrow$ float}
RMS of the non-zero couplings, $J_{\mathrm{rms}}=\sqrt{\langle
J_{ij}^2\rangle}$; accepts ndarray, DenseIsing, BQM or graph.
\end{apibox}

% ===========================================================================
\section{Annealer classes}
\label{sec:api-annealers}

Low-level API; most users should call \code{solve()}. All annealers take an
optional \code{backend="cpu"} argument and support \code{set\_seed(int)}.

\subsection{\code{Annealer}}

\begin{apibox}{Annealer(hamiltonian, schedule);\;\; run(sweeps\_per\_beta,
observer=None) $\rightarrow$ AnnealResult}
Single-chain simulated annealing (Chapter~\ref{ch:sa}). Optional
\code{MetricsObserver} / \code{StateTraceObserver}.
\end{apibox}

\subsection{\code{ReplicaAnnealer}}

\begin{apibox}{ReplicaAnnealer(hamiltonian, schedule, replicas);\;\;
run(sweeps\_per\_beta) $\rightarrow$ MultiAnnealResult}
$R$ independent SA chains, OpenMP-parallel
(Section~\ref{sec:replica-annealer}).
\end{apibox}

\subsection{\code{ParallelTemperingAnnealer}}

\begin{apibox}{ParallelTemperingAnnealer(hamiltonian, betas);\;\;
run(sweeps\_per\_step, steps, swap\_interval=1) $\rightarrow$
ParallelTemperingResult}
Classical replica exchange with the swap rule Eq.~\eqref{eq:pt-swap}
(Section~\ref{sec:classical-pt}).
\end{apibox}

\subsection{\code{SQAAnnealer}}

\begin{apibox}{SQAAnnealer(hamiltonian, schedule, trotter\_slices,
replicas=1)}
Suzuki--Trotter SQA (Chapter~\ref{ch:sqa}).
\end{apibox}

\begin{apibox}{run(sweeps\_per\_beta, worldline\_sweeps, cluster\_sweeps=0,
continuous\_time\_slices=0, observer=None) $\rightarrow$ SQAResult}
Standard $(\beta,\Gamma)$-schedule run with the moves of
Section~\ref{sec:sqa-moves}.
\end{apibox}

\begin{apibox}{run\_optimal(beta, j\_perp\_start, j\_perp\_end, eps\_tilde,
alpha=15/14, num\_steps=100, sweeps\_per\_step=20, worldline\_sweeps=0,
cluster\_sweeps=0, calib\_probes=12, calib\_sweeps=10, debug\_csv\_path="",
beta\_ramp\_fraction=0.3, beta\_ramp\_start=0.1) $\rightarrow$ SQAResult}
Adaptive $\Jp$ schedule (Chapter~\ref{ch:optimal}).
\code{eps\_tilde<=0} triggers budget calibration + inverse-CDF trajectory;
two-phase $\beta$ ramp controlled by \code{beta\_ramp\_fraction}.
Raises \code{ValueError} if \code{num\_steps==0}, if
\code{j\_perp\_end<=j\_perp\_start}, or if all sweep counts are zero.
\end{apibox}

\subsection{\code{SQAParallelTemperingAnnealer}}

\begin{apibox}{SQAParallelTemperingAnnealer(hamiltonian, betas, gammas,
trotter\_slices)}
SQA replica exchange on a $(\beta,\Gamma)$ ladder
(Chapter~\ref{ch:sqapt}).
\end{apibox}

\begin{apibox}{run(sweeps\_per\_step, worldline\_sweeps, steps,
swap\_interval=1, cluster\_sweeps=0, continuous\_time\_slices=0)
$\rightarrow$ SQAParallelTemperingResult}
Static-ladder run with the action-based swap rule
Eq.~\eqref{eq:sqapt-swap}.
\end{apibox}

\begin{apibox}{run\_optimal(num\_steps, sweeps\_per\_step, worldline\_sweeps,
eps\_tilde, alpha=15/14, j\_perp\_end=0.0, cluster\_sweeps=0,
swap\_interval=1, continuous\_time\_slices=0, calib\_probes=12,
calib\_sweeps=10, debug\_csv\_path="") $\rightarrow$
SQAParallelTemperingResult}
Shared-$\Jp$ adaptive schedule across the ladder
(Section~\ref{sec:sqapt-optimal}); same calibration semantics as the SQA
version, no two-phase ramp. \code{j\_perp\_end=0.0} is a sentinel resolving
to the coldest rung's coupling (emits a warning) --- pass an explicit value
in scripts. Raises \code{ValueError} for zero \code{num\_steps} or
\code{swap\_interval}, or if all sweep counts are zero.
\end{apibox}

\subsection{\code{CTPIMCAnnealer}}

\begin{apibox}{CTPIMCAnnealer(ising, schedule, qubits\_per\_update=1,
qubits\_per\_chain=1)\quad — general mode}
Continuous-time PIMC for arbitrary Dense/Sparse Ising
(Chapter~\ref{ch:ctpimc}).
\end{apibox}

\begin{apibox}{CTPIMCAnnealer(Lperiodic, inv\_temp\_over\_J, gamma\_over\_J,
initial\_condition=0, qubits\_per\_update=1, qubits\_per\_chain=1)\quad
— lattice mode}
Cylindrical-lattice experiments; \code{Lperiodic} must be a multiple of 6;
\code{initial\_condition}: 0 random, 1 ferro, $-1$ antiferro.
\end{apibox}

\begin{apibox}{run(sweeps\_per\_beta, reads=1) $\rightarrow$ CTPIMCResult}
Worldline cluster sweeps per schedule step; best over \code{reads} returned,
trace averaged over reads. $\Gamma$ must stay $>0$.
\end{apibox}

% ===========================================================================
\section{Result classes}
\label{sec:api-results}

\begin{center}
\footnotesize
\begin{tabular}{@{}lp{4.2cm}p{6.6cm}@{}}
\toprule
\textbf{Class} & \textbf{Produced by} & \textbf{Fields}\\
\midrule
\code{AnnealResult} & \code{Annealer.run} &
\code{best\_state}, \code{best\_energy}, \code{energy\_trace}\\
\code{MultiAnnealResult} & \code{ReplicaAnnealer.run} &
\code{global\_best\_state}, \code{global\_best\_energy}, \code{replicas},
\code{average\_energy\_trace}, \code{average\_magnetization\_trace}\\
\code{ParallelTemperingResult} & \code{ParallelTemperingAnnealer.run} &
\code{best\_state}, \code{best\_energy}, \code{final\_states},
\code{final\_energies}, \code{average\_energy\_trace},
\code{swap\_acceptance\_trace}\\
\code{SQAResult} & \code{SQAAnnealer.run/run\_optimal} &
\code{best\_state}, \code{best\_energy}, \code{energy\_trace},
\code{j\_perp\_trace}, plus calibration fields (below)\\
\code{SQAParallelTemperingResult} & \code{SQAParallelTemperingAnnealer} &
as \code{ParallelTemperingResult} + \code{j\_perp\_trace} + calibration
fields\\
\code{CTPIMCResult} & \code{CTPIMCAnnealer.run} &
\code{best\_state}, \code{best\_energy}, \code{energy\_trace}\\
\code{State} & (container) & \code{spins} ($\pm1$ int8 list),
\code{size()}\\
\bottomrule
\end{tabular}
\end{center}

\subsection{Adaptive-schedule fields on \code{SQAResult} /
\code{SQAParallelTemperingResult}}

\begin{center}
\small
\begin{tabular}{@{}lp{9.4cm}@{}}
\toprule
\textbf{Field} & \textbf{Meaning}\\
\midrule
\code{j\_perp\_trace} & Realised $\Jp$ per adaptive step; empty for
standard \code{run()}; monotone non-decreasing;
\code{len == len(energy\_trace)}.\\
\code{chi\_B\_trace} & Online $\chiB$ measured at each step.\\
\code{calibrated\_eps\_tilde} & The $\epst$ from
Eq.~\eqref{eq:calibrated-eps} (or the value you passed).\\
\code{j\_perp\_start}, \code{resolved\_j\_perp\_end} & The endpoints the run
actually used after sentinel/auto resolution.\\
\code{final\_j\_perp} & $\Jp$ reached at the last step
($\approx$ \code{resolved\_j\_perp\_end} when calibrated).\\
\bottomrule
\end{tabular}
\end{center}

% ===========================================================================
\section{Observer classes}
\label{sec:api-observers}

Observers hook into the annealing loop for diagnostics without touching the
algorithm. Attaching a sweep-level observer to \code{SQAAnnealer} disables
its OpenMP fast path (observation is inherently sequential), so use them for
analysis runs, not production sweeps.

\begin{center}
\footnotesize
\begin{tabular}{@{}lp{2.4cm}p{7.4cm}@{}}
\toprule
\textbf{Class} & \textbf{Engine} & \textbf{Captures}\\
\midrule
\code{MetricsObserver} & SA &
\code{energy\_trace}, \code{magnetization\_trace} per step;
\code{clear()} resets.\\
\code{StateTraceObserver(stride=1)} & SA &
full \code{state\_trace} plus \code{energy\_trace},
\code{step\_trace}, \code{sweep\_trace}, \code{beta\_trace} per recorded
sweep.\\
\code{SQAMetricsObserver} & SQA &
slice/replica-averaged \code{energy\_trace},
\code{magnetization\_trace}.\\
\code{SQAStateTraceObserver(stride=1)} & SQA &
\code{avg\_energy\_trace}, \code{replica\_energy\_trace},
\code{state\_trace}, \code{beta\_trace}, \code{gamma\_trace},
\code{phase\_trace} (Slice/Worldline/Cluster), \code{step\_trace},
\code{replica\_trace}, \code{sweep\_trace}; dimensions \code{replicas},
\code{slices}, \code{spins}.\\
\bottomrule
\end{tabular}
\end{center}

\section{Utility functions}

\begin{pycode}
from qanneal import magnetization, overlap

magnetization(spins)        # sum(s_i) / n
overlap(spins_a, spins_b)   # sum(a_i * b_i) / n
\end{pycode}

Both accept \code{list}, \code{np.ndarray}, or a \code{State}.
Magnetisation tracks global ordering during an anneal; overlap between two
reads' solutions measures how reproducibly a problem's ground state is
found (overlap $\approx1$: same solution up to global flip).
